\chapter{Migrating from other TUI libraries}
\label{ch:migrating}

\begin{aside}
You have an app in ncurses, ftxui, blessed, ratatui. Moving
to \maya{} is mostly a port, mostly mechanical. The patterns
translate.
\end{aside}

\section{ncurses}

ncurses is C, imperative. \maya{} is C++, declarative. The
big change is moving from ``write to screen at coordinates''
to ``return an element tree''.

A typical ncurses loop:

\begin{lstlisting}
mvprintw(y, x, "%s", text);
refresh();
ch = getch();
// handle ch
\end{lstlisting}

Becomes:

\begin{lstlisting}
return text(s) | position(x, y);
\end{lstlisting}

Layout calculates positions; you don't.

\section{ftxui}

ftxui has a similar element model. Direct port: element
names mostly match. \maya{}'s signals replace ftxui's
\code{Component} state.

\section{blessed (Node.js)}

blessed is event-driven and component-based. \maya{}'s
event handlers and signals map cleanly.

\section{ratatui (Rust)}

Ratatui draws into a Frame imperatively each render. \maya{}
uses a tree. The conversion is moving each Frame draw call
into an element.

\section{The data layer}

Whatever your app's data layer is, \maya{} doesn't care. Put
it behind signals and consume from views.

\section{Style and theme}

Each library has its own colour model. Map your theme to
\maya{}'s palette type. Hex literals work universally.

\section{Testing}

Snapshots port directly: render your view, compare to
golden. Adjust the format if needed.

\section{Recap}

\begin{recap}
\begin{itemize}[leftmargin=*]
  \item ncurses: write-to-screen \textrightarrow{}
    return-an-element.
  \item ftxui, ratatui, blessed: mostly mechanical.
  \item Data layer is independent.
  \item Snapshots port directly.
\end{itemize}
\end{recap}

\section*{Workshop}

\begin{enumerate}[leftmargin=*]
  \item Port one screen from your old library to \maya{}.
  \item Compare line counts and clarity.
  \item Identify any \maya{} feature missing from the old
    one.
\end{enumerate}
